\BourbakiExerciseGroup

\begin{enumerate}
\def\labelenumi{\arabic{enumi})}
\tightlist
\item
  \begin{enumerate}
  \def\labelenumii{\alph{enumii})}
  \tightlist
  \item
    Show that the ring \(\mathbf{Z}_p\) of \(p\) -adic integers {[}§ 6, Exercise 23 a){]} is isomorphic to the inverse limit of the sequence of discrete rings \(\mathbf{Z}/(p^n)\) , \(f_{nm}: \mathbf{Z}/(p^m) \to \mathbf{Z}/(p^n)\) for \(n \leqslant m\) being the canonical homomorphism when \(\mathbf{Z}/(p^n)\) is considered as a quotient ring of \(\mathbf{Z}/(p^m)\) .
  \end{enumerate}
\end{enumerate}

\begin{enumerate}
\def\labelenumi{\alph{enumi})}
\setcounter{enumi}{1}
\item
  For each integer \(n > 0\) , let \(G_n\) be the group \(Z\) of rational integers, and let \(X_n\) be its quotient group \(Z/(p^n)\) ( \(p\) prime), \(G_n\) and \(X_n\) carrying the discrete topology; consider \((G_n)\) as an inverse system, \(G_m \to G_n\) being the identity mapping; and consider \((X_n)\) as an inverse system, \(f_{nm}: X_m \to X_n\) being the canonical homomorphism (see a) for \(n \leqslant m\) . Then the orbit of each point of \(X_n\) with respect to \(G_n\) is compact, but the orbit of a point \(x = (x_n)\) of \(X = \varprojlim X_n\) with respect to \(G = \varprojlim G_n\) is not compact, and is not isomorphic to the inverse limit of the orbits of the \(x_n\) ; also the canonical mapping \(X/G \to \varprojlim X_n/G_n\) is not injective. c) For each \(n > 0\) , let \(G_n\) be the subgroup \(p^n Z\) of \(Z\) ( \(p\) prime), and let \(X_n\) be the group \(Z\) , where \(G_n\) and \(X_n\) carry the discrete topology and \(G_n\) operates by translations on \(X_n\) . Consider \((G_n)\) and \((X_n)\) as inverse systems, \(X_m \to X_n\) being the identity mapping and \(G_m \to G_n\) the canonical injection ( \(n \leqslant m\) ). Then the stabilizer of each point of \(X_n\) (with respect to \(G_n\) ) is compact, the orbit of each point of \(X = \varprojlim X_n\) with respect to \(G = \varprojlim G_n\) is compact, but the canonical mapping \(X/G = \varprojlim X_n/G_n\) is not surjective.
\item
  Construct, using b) and c), an example of an inverse system of spaces with operators \((\mathbf{X}_{\alpha})\) with respect to an inverse system of groups \((\mathbf{G}_{\alpha})\) , such that, if \(X = \varprojlim X_{\alpha}\) and \(G = \varprojlim G_{\alpha}\) , the canonical mapping \(X/G \to \varprojlim X_{\alpha}/G_{\alpha}\) is neither injective nor surjective.
\end{enumerate}

\begin{enumerate}
\def\labelenumi{\arabic{enumi})}
\setcounter{enumi}{1}
\tightlist
\item
  Let \((\mathbf{X}_{\alpha}, f_{\alpha \beta})\) be an inverse system of non-empty sets such that \(\varprojlim \mathbf{X}_{\alpha} = \emptyset\) . Let \(\mathbf{G}_{\alpha}\) denote the free \(\mathbf{Z}\) -module of formal linear combinations of elements of \(\mathbf{X}_{\alpha}\) , with coefficients in \(\mathbf{Z}\) , and let \(g_{\alpha \beta}\) denote the homomorphism \(\mathbf{G}_{\beta} \to \mathbf{G}_{\alpha}\) which agrees with \(f_{\alpha \beta}\) on \(\mathbf{X}_{\beta}\) . Show that the group \(\mathbf{G} = \varprojlim \mathbf{G}_{\alpha}\) consists of the identity element alone. {[}Argue by contradiction: if \(z = (z_{\alpha})\) is an element of \(\varprojlim \mathbf{G}_{\alpha}\) , consider, for each \(\alpha\) , the finite set \(\mathbf{F}_{\alpha}\) of elements of \(\mathbf{X}_{\alpha}\) whose coefficient in \(z_{\alpha}\) is not zero, and observe that \(f_{\alpha \beta}(\mathbf{F}_{\beta}) = \mathbf{F}_{\alpha}\) .{]} Hence construct an example of an inverse system \((\mathbf{G}_{\alpha}, g_{\alpha \beta})\) of groups such that the \(g_{\alpha \beta}\) are surjective, the \(\mathbf{G}_{\alpha}\) are infinite and \(\varprojlim \mathbf{G}_{\alpha}\) consists only of the identity element.
\end{enumerate}

¶ 3) a) Show that every totally disconnected compact group is the inverse limit of a family of finite discrete groups (use Exercise 18 of § 4). b) Let G be a totally disconnected compact group and let L be a closed subgroup of G. Show that there is a continuous section G/L → G associated with the canonical mapping G → G/L. {[}Use a) and Proposition 3 of § 3; for each α consider the finite set Fα of sections Gα/Lα → Gα, and remark that these sets from an inverse system with respect to the canonical surjective mappings hαβ : Fβ → Fα{]}.

\begin{enumerate}
\def\labelenumi{\arabic{enumi})}
\setcounter{enumi}{3}
\item
  Let \(G\) be a Hausdorff topological group and let \((H_{\alpha})\) be a family of compact normal subgroups of \(G\) , directed with respect to the relation \(\supset\) , and satisfying condition (AP) of no. 3. Let \((L_{\alpha})\) be a family of closed subgroups of \(G\) such that \(H_{\alpha} \subset L_{\alpha}\) for each \(\alpha \in I\) and \(L_{\alpha} = H_{\alpha}L_{\beta}\) for \(\alpha \leqslant \beta\) , and let \(L = \bigcap_{\alpha} L_{\alpha}\) . Show that \(L_{\alpha} = H_{\alpha}L\) for each \(\alpha\) (use Proposition 3).
\item
  Let G be a Hausdorff topological group and let X be a Hausdorff topological space on which G operates continuously; let \((\mathrm{H}_{\alpha})\) be a family of compact normal subgroups of G, directed with respect to the relation \(\supset\) , and satisfying condition (AP) of no. 3. Let \(X_{\alpha}\) denote \(X/H_{\alpha}\) . Show that the canonical mapping \(X \to \varprojlim X_{\alpha}\) is a homeomorphism.
\end{enumerate}

¶ * 6) Let \((\mathbf{G}_n, f_{nm})_{n \in \mathbb{N}}\) be the inverse system of compact groups such that \(\mathbf{G}_n = \mathbf{T} = \mathbf{R} / \mathbf{Z}\) for each \(n\) , and \(f_{nm}\) is the continuous homomorphism \(x \to p^{m-n}x\) of \(\mathbf{T}\) onto itself, for \(n \leqslant m\) ( \(p\) being a given prime number). The topological group \(\mathbf{T}_p = \varprojlim \mathbf{G}_n\) is called the \(p\) -adic solenoid; it is a compact connected commutative topological group.

\begin{enumerate}
\def\labelenumi{\alph{enumi})}
\item
  For each n, the continuous homomorphism \(f_{n}: T_{p} \to G_{n}\) is surjective, and its kernel is isomorphic to the group \(Z_{p}\) of p-adic integers {[}cf.~Exercise 1 a){]}.
\item
  Let \(\varphi\) be the canonical homomorphism \(\mathbf{R} \to \mathbf{R} / \mathbf{Z} = \mathbf{T}\) . For each \(x \in \mathbf{R}\) , put \(\theta(x) = (\varphi(x / p^n))_{n \in \mathbb{N}}\) ; show that \(\theta\) is an injective continuous homomorphism of \(\mathbf{R}\) into \(\mathbf{T}_p\) , and that \(\theta(\mathbf{R})\) is a dense subgroup of \(\mathbf{T}_p\) .
\item
  Let \(\mathbf{I}\) be an open interval in \(\mathbf{R}\) , with centre \(0\) and length \(< 1\) . Show that the subspace \(\overline{f}_{0}^{1}(\varphi(\mathbf{I}))\) of \(\mathbf{T}_p\) is homeomorphic to the product \(\mathbf{I} \times \mathbf{Z}_p\) . In particular, the group \(\mathbf{T}_p\) is not locally connected.
\item
  Show that every closed subgroup \(\mathbf{H}\) of \(\mathbf{T}_p\) , other than \(\mathbf{T}_p\) or \(\{0\}\) , is totally disconnected and isomorphic to a group of the form \(\mathbf{Z}/n\mathbf{Z}\) or \((\mathbf{Z}/n\mathbf{Z}) \times \mathbf{Z}_p\) , where \(n\) is an integer prime to \(p\) (use Proposition 3 of no. 3).
\item
  Show that \(\mathbf{T}_p\) is an indecomposable compact connected space, i.e.~that \(\mathbf{T}_p\) cannot be covered by two compact connected sets \(\mathbf{P}\) , \(\mathbf{Q}\) neither of which is equal to \(\mathbf{T}_p\) . (Observe that there is an integer \(n\) such that \(f_n(\mathbf{P}) \neq \mathbf{G}_n\) and \(f_n(\mathbf{Q}) \neq \mathbf{G}_n\) , and examine the sets \(f_{n+1}(\mathbf{P})\) and \(f_{n+1}(\mathbf{Q})\) to get a contradiction). *
\end{enumerate}

